% -------------------------------
% Plantilla del TFG en la ESIIAB-UCLM. Versión 0.0 (Preliminar)
%
% Luis de la Ossa
% 
% Compilar con XeLaTeX 
% -------------------------------


% -------------------------------
% Segunda portada
% -------------------------------
\cleardoublepage
\setcounter{page}{1} \null

\begin{textblock*}{3cm}(3cm,2cm) 
\begin{flushleft}
\includegraphics[height=2cm]{figs/logouclm.png}
\end{flushleft}
\end{textblock*}


\begin{textblock*}{2.5cm}(15.95cm,2.45cm) 
\begin{flushright}
\includegraphics[height=1.4cm]{figs/euroinf.jpg}
\end{flushright}
\end{textblock*}

\begin{textblock*}{4cm}(13.5cm,2cm)  % Ancho ajustado a 4cm para evitar overfull hbox
\begin{flushright}
\includegraphics[height=2cm,keepaspectratio]{figs/logoesii.png}
\end{flushright}
\end{textblock*}



\begin{textblock*}{\textwidth}(3cm,8cm) 
\begin{center}\doublespacing 
{\fontsize{14pt}{3pt}\selectfont \bft{TRABAJO DE PROCESOS DE\\ INGENIERÍA DEL SOFTWARE}}

\medskip
{\fontsize{18pt}{4pt}\selectfont Grado en Ingeniería Informática}

\medskip
{\fontsize{16pt}{4pt}\selectfont \espec}
\end{center}
\end{textblock*}


\begin{textblock*}{\textwidth}(3cm,14cm) 
\begin{center}\setstretch{2}
{\fontsize{22pt}{4pt}\selectfont \bft{\titulo}}
\end{center}
\end{textblock*}


\begin{textblock*}{\textwidth}(3cm,20.5cm) 
\begin{flushleft}\doublespacing
{\fontsize{14pt}{4pt}\selectfont \bft{Autor:} \autor}

\ifdefined\director{\fontsize{14pt}{4pt}\selectfont \bft{Tutor:} \director}\fi

\ifdefined\codirector{\fontsize{14pt}{4pt}\selectfont \bft{Co-Tutor:} \codirector}\fi
\end{flushleft}
\end{textblock*}


\begin{textblock*}{\linewidth}(3cm,25cm) 
\begin{flushright}
{\fontsize{14pt}{4pt}\selectfont \fecha}
\end{flushright}
\end{textblock*}


% -------------------------------
% Declaración de autoría
% -------------------------------
\cleardoublepage
\thispagestyle{plain}
\begin{center}
\Large{\bft{Declaración de autoría}}
\end{center}
\vskip1cm

Yo, \autor, declaro que soy el autor del trabajo titulado ``\titulo'', que el citado trabajo no infringe las leyes en vigor sobre propiedad intelectual, y que todo el material no original contenido en dicho trabajo está apropiadamente atribuido a sus legítimos autores, detallándose asimismo el uso responsable de la inteligencia artificial en el Anexo~\ref{anx:ia}.

\vspace*{2cm}
\begin{center}
Albacete, \fecha

\vskip3cm

Fdo.: \autor
\end{center}


% -------------------------------
% Resumen
% -------------------------------
\cleardoublepage
\thispagestyle{plain}
\begin{center}
\Large{\bft{Resumen}}
\end{center}
\vskip1cm

Esta memoria describe el diseño, la implementación y el proceso de desarrollo completo del proyecto \textit{La Grada Colchonera}, una aplicación SaaS organizada en capas con gestión de usuarios y un juego web de tipo \textit{survivor} en pixel art, desarrollado como parte de la asignatura Procesos de Ingeniería del Software (curso 2026-2027). El trabajo abarca dos sprints: el Sprint 1 (29/09–27/10/2026, Hitos 1–4) construye la base SaaS con cuentas, sesión, roles, PostgreSQL, correo, Google OAuth, errores, logs, CI/CD y despliegue en Google Cloud Run; el Sprint 2 (23/10–29/10/2026) implementa la lógica del juego: motor survivor con oleadas, jefes, armas automáticas, mejoras, mundo procedural, club persistente en servidor con validación anti-trampas, logros con historia del Atlético de Madrid, música chiptune y sonidos de grada. Se detalla exhaustivamente la arquitectura en capas, el modelo de datos, la gestión segura de sesiones y contraseñas mediante bcrypt, la metodología Scrum en DoneTonic, GitHub Flow, trazabilidad requisito → PBI → tarea → PR, y la evaluación con 276 pruebas automatizadas y revisión frase a frase contra el enunciado.

\medskip
\noindent\textbf{Palabras clave:} SaaS, Scrum, GitHub Flow, Node.js, PostgreSQL, Google Cloud Run, survivor, pixel art, anti-cheat, Ingeniería del Software.

% -------------------------------
% Abstract
% -------------------------------
\cleardoublepage
\thispagestyle{plain}
\begin{center}
\Large{\bft{Abstract}}
\end{center}
\vskip1cm

This memoir describes the design, implementation, and development process of the \textit{La Grada Colchonera} project, a layered SaaS application with user management and a web \textit{survivor}-style game in pixel art, developed for the Software Engineering Processes course (2026-2027). The work spans two sprints: Sprint 1 (29/09–27/10/2026, Milestones 1–4) builds the SaaS foundation with accounts, sessions, roles, PostgreSQL, email, Google OAuth, error handling, logging, CI/CD, and deployment on Google Cloud Run; Sprint 2 (23/10–29/10/2026) implements the game logic: survivor engine with waves, bosses, automatic weapons, upgrades, procedural world, server-persisted club with anti-cheat validation, achievements tied to Atlético de Madrid history, chiptune music, and crowd sounds. The memoir thoroughly details the layered architecture, data model, secure session and password management with bcrypt, Scrum methodology in DoneTonic, GitHub Flow, requirement → PBI → task → PR traceability, and evaluation with 276 automated tests and phrase-by-phrase review against the specification.

\medskip
\noindent\textbf{Keywords:} SaaS, Scrum, GitHub Flow, Node.js, PostgreSQL, Google Cloud Run, survivor, pixel art, anti-cheat, Software Engineering.
